% Source-derived chapter 02: Spherical varieties and their arc spaces
% Original source: intersection-complexes-unramified-l-factors
\chapter{Spherical varieties and their arc spaces}
\label{intersection-complexes-unramified-l-factors-chapter-02}
\source{intersection-complexes-unramified-l-factors}

\begin{remark}[Source section]
\label{intersection-complexes-unramified-l-factors-chapter-02-source}
\source{intersection-complexes-unramified-l-factors}
\source{intersection-complexes-unramified-l-factors}
This chapter reproduces the corresponding section of the retrieved
LaTeX source, including its definitions, statements, examples, and
proofs. It has not yet been translated into Lean.
\notready
\end{remark}

% The source section title was: Spherical varieties and their arc spaces
\label{section-spherical}
\source{intersection-complexes-unramified-l-factors}

\subsection{Spherical varieties}
\label{sect:spherical}
\source{intersection-complexes-unramified-l-factors}

A spherical variety over $k$ is a normal variety with an open $B$-orbit. 
Let $X$ be an affine $G$-spherical variety over $k$. Let $X^\circ$ denote the open $B$-orbit. We choose and fix a point $x_0 \in X^\circ(k)$ and let $H$ denote its stabilizer.
Let $X^\bullet= H\bs G $ denote the open $G$-orbit. 

The quotient $X^\circ\sslash N$ has an action of the universal Cartan $T=B/N$, and is a torsor for a quotient torus $T\twoheadrightarrow T_X$. By our choice of base point, we can identify this torsor with $T_X$. In the rest of this paper, we will assume that our spherical variety satisfies $T_X=T$; however, for now we proceed with general definitions.

All important combinatorial invariants of the spherical variety live in the rational vector space $\mathfrak t_X^*$ spanned by the character group $\Lambda_X$ of this torus, or in the dual vector space $\mathfrak t_X$, containing the dual lattice $\ch\Lambda_X$. By ``lattice points'', below, we will mean points belonging to one of these lattices. The spaces $\mathfrak t_X, \mathfrak t_X^*$ are the root and coroot space for the dual group $\check G_X$ of $X$.\footnote{The \emph{Gaitsgory--Nadler dual group} was defined in a Tannakian way in \cite{GN}, but not completely identified in all cases. A combinatorial description  of a dual group (presumably the same) was consequently afforded by Knop and Schalke \cite{Knop-Schalke}. The invariants that we present here are those of Knop and Schalke, which match standard invariants of the theory of spherical varieties. For this paper, however, this distinction between constructions of the dual group is immaterial, as we impose the condition that $T_X= T$ and ``all simple roots of $G$ are spherical roots of type $T$'', which implies that in both versions of the dual group, $\check G_X= \check G$.} The antidominant Weyl chamber for $\check G_X$ in $\mathfrak t_X$ is denoted by $\Cal V$ in the theory of spherical varieties, because it coincides with the so-called cone of $G$-invariant valuations, see \cite{KnLV}. Up to this point, all data depend only on the open $G$-orbit $X^\bullet$, not on its affine embedding $X$.

The affine embedding $X$ of $X^\circ$ defines an affine toric embedding $X\sslash N$ of $T_X$, described by the cone $\Cal C_0(X) \subset \mathfrak t_X$ whose lattice points are all cocharacters $\check\lambda$ into $T_X$ such that $\lim_{t\to 0} t^{\ch\lambda}\in X\sslash N$. We will denote by $\mathfrak c_X$ the monoid of lattice points $\check\Lambda_X \cap \Cal C_0(X) \subset \mathfrak t_X$, and by $\mathfrak c_X^-$ its intersection with the cone $\Cal V$ of invariant valuations. 
The cone $\Cal C_0(X)\subset \mathfrak t_X$ has a canonical set of generators $\ch\nu_D$, the valuations associated to all $B$-stable divisors $D\subset X$. The set of all irreducible $B$-stable divisors in $X$ will be denoted by $\mathcal D(X)$, and by ``valuation associated'' we mean the restriction of the corresponding valuation to the group of nonzero $B$-eigenfunctions: $\check \nu_D: k(X)^{(B)} \to \mathbb Z$, which factors through the character group $\Lambda_X$ and hence can be identified with an element of $\check\Lambda_X\subset \mathfrak t_X$. The map $\mathcal D(X)\ni D\mapsto \check\nu_D\in \check\Lambda_X$ will be denoted by $\varrho_X$. Inside of $\mathcal D(X)$ there is a distinguished subset $\mathcal D$, depending only on the open $G$-orbit $X^\bullet$, which consists of the closures of $B$-stable divisors in $X^\bullet$; those are called \emph{colors}. We will often abuse language and write ``colors'' for the images of $\mathcal D$ in $\ch\Lambda_X$.





We remark that the valuation map $\varrho_X$ may fail to be injective, but this can only happen when two colors have the same image. If this is the case for $X^\bullet = H\backslash G$, there is always a torus covering of it such that all colors have distinct valuations (see \S\ref{sect:freemonoid}); for example, $\GL_1\backslash \PGL_2$ has two colors with valuation $\frac{\check\alpha}{2}$, but their preimages in $\GL_1\backslash\GL_2$ (where $\GL_1$ is embedded as the general linear group of a one-dimensional subspace) induce different valuations. 
In any case, colors give rise to a map $\mathbb N^{\mathcal D}\to \mathfrak c_X$, whose image we denote by $\mathfrak c_X^{\Cal D}$. 

We define an ordering $\preceq$ on $\ch\Lambda_X$ by postulating that $\ch\lambda \preceq \ch\lambda'$ if 
$\ch\lambda'-\ch\lambda$ can be written as a non-negative \emph{integral} combination of the valuations $\ch\nu_D$, with $D\in \mathcal D$, i.e., if $\ch\lambda'-\ch\lambda \in \mathfrak c_X^{\Cal D}$. We use the symbol $\le $ for the ordering on $\ch\Lambda_G$ defined by the positive coroots of $G$, that is, $\ch\lambda \le \ch\lambda'$ iff $\ch\lambda'-\ch\lambda$ is a sum of positive coroots. Notice that the ordering $\preceq$ is not defined simply in terms of the cone spanned by the $\ch\nu_D$'s: there can be non-comparable lattice points in this cone (and similarly for the ordering $\le$). 
As we will see later, this ordering describes the closure relations on the global model 
of the arc space of $X$.

\subsubsection{Spherical roots of type $T$} 
\label{subsection:typeT}
\source{intersection-complexes-unramified-l-factors}

From Section \ref{sect:models} onwards we assume that $T_X=T$, equivalently, $B$ acts simply transitively on $X^\circ$. From Section \ref{sect:Heckeact} on we assume, further, that \emph{all simple roots of $G$ are spherical roots of type $T$}. Let us explain what this means: 
For a simple root $\alpha$ of $G$, let $P_\alpha \supset B$ denote 
the corresponding parabolic of semisimple rank one. The quotient $P_\alpha/\mathfrak R(P_\alpha)$ by its radical is isomorphic to $\PGL_2$, and the invariant-theoretic (or geometric) quotient $X^\circ P_\alpha/\mathfrak R(P_\alpha)$ is a spherical variety for $\PGL_2$. In characteristic zero, over an algebraically closed field, those belong to one of the following types, see \cite[Lemma 3.2]{KnBorbit}:
\begin{itemize}
 \item a point: $\PGL_2\backslash\PGL_2$;
 \item type $T$: $\mathbb G_m\backslash \PGL_2$;
 \item type $N$: $\mathcal N(\mathbb G_m)\backslash\PGL_2$, where $\mathcal N$ denotes normalizer;
 \item type $U$: $S\backslash \PGL_2$, where $N\subset S\subset B$. 
\end{itemize}

In positive characteristic there are some more cases, investigated by Knop in \cite{Knop-sphericalroots}. 

Our assumption from Section~\ref{sect:Heckeact} onwards is that for every simple root $\alpha$, this $\PGL_2$-spherical variety is isomorphic to $\mathbb G_m\backslash \PGL_2$. Our assumptions imply that the stabilizer in $P_\alpha$ of a point on the open orbit is isomorphic to $\mathbb G_m$, and that there are precisely two colors $D_\alpha^+, D_\alpha^-$ contained in $X^\circ P_\alpha$ (the $\pm $ labeling is arbitrary). We will denote the set of these two elements by $\mathcal D(\alpha)\subset\mathcal D$; notice that these sets are not disjoint as $\alpha$ varies. Moreover, the associated valuations satisfy (see \cite[\S 3.4]{Lun97}):
\[ \ch\nu_{D^+_\alpha} = -s_\alpha \ch\nu_{D^-_\alpha},\]
\[ \ch\nu_{D^+_\alpha} + \ch\nu_{D^-_\alpha} = \ch\alpha, \]
and finally an element $D\in \mathcal D$ belongs to $\mathcal D(\alpha)$ iff $\left<\alpha,\ch\nu_{D}\right> >0$ (in which case $\left<\alpha,\ch\nu_{D}\right> =1$, by the above). 

\begin{rem}
\source{intersection-complexes-unramified-l-factors}
\notready
 Our assumptions above are over the algebraically closed field $k$. Over a finite field $\mathbb F$, the Galois group acts on the set of colors, compatibly with its action on the set of simple roots. If, for example, $G$ is split, each set $\mathcal D(\alpha)$ is preserved by Frobenius, and the stabilizer of a point on the open $P_\alpha$-orbit is a form of $\mathbb G_m$. If that form is split, Frobenius acts trivially on $\mathcal D(\alpha)$; if not, it permutes the two colors.
\end{rem}


\begin{rem}
\source{intersection-complexes-unramified-l-factors}
\notready\label{rem:valuation}
\source{intersection-complexes-unramified-l-factors}
A very straightforward way to compute the valuations $\ch\nu_{D^{\pm}_\alpha}$ in any example is the following: If all simple roots are spherical roots of type $T$ and $T_X=T$, the stabilizer $S$ of a point in the open $P_\alpha$-orbit on $X$ is a subgroup isomorphic (over the algebraic closure) to $\mathbb G_m$. Choose a Borel subgroup $B\subset P_\alpha$ containing $S$. An isomorphism $\mathbb G_m\simeq S$ gives rise to a cocharacter $\ch\nu:\mathbb G_m \to B\to T$, and we can choose this isomorphism so that $\left<\alpha, \ch\nu \right> = 1$. There are two choices for $B$, and they correspond to the valuations $\ch\nu_{D^{\pm}_\alpha}$.
\end{rem}





\subsection{Affine degeneration, and assumptions in positive characteristic} \label{sect:affine-degeneration}  \label{assumptions-finitefield}
\source{intersection-complexes-unramified-l-factors}

In characteristic zero, there is a well-known affine family degenerating $X$ to a horospherical variety \cite{Pop86}, \cite[\S 5.1]{GN}, which in turn is related to its degeneration to the normal bundle of orbits in a smooth toroidal (e.g., a ``wonderful'') compactification of $X$, 
\cite{Brion07}, \cite[\S 2.5]{SV}. These constructions sometimes fail in positive characteristic, and therefore \textbf{the statements of this subsection should be considered as assumptions in positive characteristic}. These assumptions will be imposed on $X$ for the remainder of the paper. Notice that the Luna--Vust theory of spherical embeddings holds in arbitrary characteristic over an algebraically closed field by \cite{KnLV}. 


\subsubsection{}
We define a filtration $\Cal F_{\lambda}$ on $k[X]$ for $\lambda \in \Lambda_X$ 
by letting $\Cal F_{\lambda}$ consist of all $f \in k[X]$
such that every highest weight $\mu \in \Lambda_X$ of the rational $G$-module 
generated by $f$ satisfies $\brac{\lambda-\mu, \Cal V} \le 0$.
The affine degeneration $\mathscr X$ is the affine variety defined to be the
spectrum of the Rees algebra associated to the filtration above: 
\[ k[\mathscr X] = \bigoplus_{\lambda\in \Lambda_X} \Cal F_\lambda \ot e^\lambda \subset k[X \times T_X] 
\] 
where $e^\lambda \in k[T_X]$ denotes the character corresponding to $\lambda \in \Lambda_X$. This family is naturally equipped with an action of the product $G\times T_X$.

Note that $k[\mathscr X]$ contains $\bigoplus_{\brac{\lambda,\Cal V}\le 0} k e^\lambda$. 
Define $\ol{T_{X,\mss}}$ to be the spectrum of $\bigoplus_{\brac{\lambda,\Cal V}\le 0} k e^\lambda$. 
This is an affine toric variety with open orbit isomorphic to the
quotient $T_{X,\mss}$ of $T_X$ by the subtorus $\mathcal Z(X)^0$ (the ``connected center of $X$'') generated by cocharacters in $\Cal V \cap (-\Cal V)$.

In summary, we get a $G \times T_X$-equivariant map 
\[ \mathscr X \to \ol{T_{X,\mss}}, \] 
which we consider as our affine family of degenerations.

\subsubsection{} 
For $a \in \ol{T_{X,\mss}}(k)$, let $X_a$ denote the fiber of 
$\mathscr X \to \ol{T_{X,\mss}}$ over $a$. 
From the definition of the filtration $\Cal F_\lambda$, it follows that 
$k[X_a]^N = k[X]^N$. 
The following ``multiplicity-free'' characterization of spherical varieties (in arbitrary
characteristic) implies that each fiber $X_a$ is spherical.

\begin{thm}[{\cite{VinbergKimelfeld}, \cite[Theorem 25.1]{Timashev}}]
\source{intersection-complexes-unramified-l-factors}
\notready
A normal quasi-affine variety $X$ is spherical if and only if 
the non-zero weight spaces of $k[X]^N$ are all $1$-dimensional.
\end{thm}

Since the torus $T_X$ is determined by $k[X]^N$, we have a family of spherical varieties
$X_a$ with associated torus $T_{X_a} \cong T_X$. 

We also have $k[\mathscr X]^N = \bigoplus_{\brac{\lambda-\mu,\Cal V}\le 0} k[X]^{(B)}_\mu \ot e^\lambda$ where $k[X]^{(B)}_\mu$ is the $1$-dimensional $B$-eigenspace of weight $\mu$. 
This gives a canonical isomorphism 
\begin{equation} \label{e:scrXsslashN} 
\source{intersection-complexes-unramified-l-factors}
    \mathscr X/\!\!/N \cong X/\!\!/N \times \ol{T_{X,\mss}}. 
\end{equation}
Define $\mathscr X^\circ$ to be the preimage of 
$T_X \times \ol{T_{X,\mss}} \subset \mathscr X/\!\!/N$ under the projection
$\mathscr X \to \mathscr X/\!\!/N$. 
Equivalently, $\mathscr X^\circ$ is the union of the open $B$-orbits of the fibers $X_a$. 



It will be convenient to lift the isomorphism \eqref{e:scrXsslashN} to $\mathscr X^\circ$.
\begin{prop}
\source{intersection-complexes-unramified-l-factors}
\notready \label{prop:noncanon-section}
\source{intersection-complexes-unramified-l-factors}
There is a (non-canonical) $B$-equivariant isomorphism $\mathscr X^\circ \cong X^\circ \times \ol{T_{X,\mss}}$ over $\ol{T_{X,\mss}}$, compatible with \eqref{e:scrXsslashN}.
\end{prop}

We will comment on the proof of this proposition in conjunction with Theorem \ref{thm:localstructure} below. 

The affine degeneration is closely related to compactifications of the spherical variety. 
Namely, let $\mathscr X^\bullet \subset \mathscr X$ denote the union of the $G$-translates of $\mathscr X^\circ$. (It is the open subvariety which
specializes over each fiber $X_a$ to the open $G$-orbit.)
The quotient $\mathscr X^\bullet/ T_X$ turns out to be a proper embedding of $X^\bullet/\mathcal Z(X)^0$. For applications, one is interested in proper embeddings of $X^\bullet$ itself, and preferably smooth ones. Without getting into the details of the construction (we point the reader to \cite{KnLV}), we formulate the following result on the existence and local structure of ``smooth toroidal compactifications'':

Define $P(X) \supset B$ to be the parabolic subgroup of $G$ equal to 
$\{ g \in G \mid X^\circ \cdot g = X^\circ \}$, and let $N_{P(X)}$ denote its unipotent radical. Recall that we have fixed a base point $x_0\in X^\circ$, for convenience.

\begin{thm}
\source{intersection-complexes-unramified-l-factors}
\notready\label{thm:localstructure}
\source{intersection-complexes-unramified-l-factors}
 There is a proper, smooth, $G$-equivariant embedding $X^\bullet \hookrightarrow \ol X$, a Levi\footnote{Of course, there is no distinguished Levi in $P(X)$ abstractly, but the choice of base point $x_0$ sometimes imposes restrictions on the Levi subgroups that work; for example, the derived subgroup of $L$ must stabilize $x_0$. In any case, the assumptions on $X$ in the main body of this paper, that $B$ acts simply transitively on $X^\circ$ and $\Cal V=$ the antidominant cone, imply that any $L\subset P(X)= B$ satisfies the Local Structure Theorem.} $L \subset P(X) \subset G$,  and a smooth toric embedding $\overline{T_X}$ of the quotient $T_X$ of $L$, such that the action map 
 \[ L \times N_{P(X)} \ni (l, u) \mapsto x_0 lu\]
 descends to $T_X$ and extends to an open embedding
 \[ \overline{T_X} \times N_{P(X)} \hookrightarrow \ol X,\]
 whose image is the union of all open $B$-orbits on $\ol X$. Moreover, the support of the fan describing the toric variety $\overline{T_X}$ is the cone $\Cal V$ of invariant valuations.
\end{thm}

\begin{rem}
\source{intersection-complexes-unramified-l-factors}
\notready
 Note that the embedding that we denoted above by $\ol{T_{X,\mss}}$ is associated to the image of the \emph{opposite} cone $-\Cal V$ modulo $\Cal V\cap (-\Cal V)$. Thus, there is a map from $\overline{T_X}$ to $\ol{T_{X,\mss}}$ only after inversion in $T_X$. We allow ourselves this notational flaw, since $\overline{T_X}$ will not be used beyond this subsection and the next.
\end{rem}


Theorem \ref{thm:localstructure} is the local structure theorem of Brion--Luna--Vust \cite[Th\'eor\`eme 3.5]{BLV}, applied to smooth toroidal compactifications \cite{KnLV}. We outline its proof due to Knop \cite[Theorem 2.3]{KnMotion}, which also applies to Proposition \ref{prop:noncanon-section}. The main issue is how to choose the pair $(x_0,L)$ appropriately. 
Let $\Cal B$ denote the flag variety of $G$, and assume for the moment that no choice of $B$, $x_0$ has been made.
One considers triples $(x,B,\chi) \in X\times\mathcal B\times \mf t_X^*$ such that 
$x$ lies in the open $B$-orbit. Out of these data one constructs elements in the cotangent bundle $T^*X^\bullet$ as follows: If $\chi$ is the differential of a character (also to be denoted by $\chi\in \Lambda_X$), the corresponding cotangent vector is the logarithmic differential $\left.\frac{d f_\chi}{f_\chi}\right|_x$ at $x$ of a rational $B$-eigenfunction with eigencharacter $\chi$, and for general $\chi$ we extend this construction by linearity. This gives rise to a map 
\[ (X\times\mathcal B)^\bullet \times \mf t_X^* \to T^*X^\bullet \]
(where the bullet denotes ``open $G$-orbit'') with dense image. We can also apply this construction to the family $\mathscr X^\bullet$, obtaining vectors in the \emph{relative} cotangent bundle over $\ol{T_{X,\mss}}$. Composing with the moment map $T^*X^\bullet \to \mathfrak g^*$, Knop shows that if we choose a sufficiently generic, semisimple vector $\xi$ in the image, and take $L=\on{Stab}_G(\xi)$ and $x_0$ a point of a triple $(x_0, B, \chi)$ in its fiber, Theorem \ref{thm:localstructure} is satisfied. The same argument applies to prove Proposition \ref{prop:noncanon-section}, as follows: Instead of fixing $x_0$, fix first just the Borel subgroup $B$, and consider all triples $(x,B,\chi)$ over $\xi$, now with $x \in \mathscr X^\circ$ (defined with respect to this Borel). The set of these triples is a 
$T_X$-equivariant section of the map 
$\mathscr X \to \mathscr X\sslash N \cong X\sslash N \times \ol{T_{X,\mss}}$ 
over $T_X\times \ol{T_{X,\mss}}$, where $T_X$ acts diagonally on the latter. 



\subsection{The formal loop space}
\label{sect:formalloop}
\source{intersection-complexes-unramified-l-factors}
For any $k$-scheme $X$, define the space of formal arcs by
\[ \msf L^+ X(R) = X(R\tbrac t) \] 
for a test ring $R$. 
It is well-known (cf.~\cite[Proposition 1.2.1]{KV}) that $\msf L^+ X$ is representable by a scheme (of infinite type), which is
the projective limit of the schemes $\msf L^+_n X,\, n \in \mbb N$, 
representing the spaces of $n$-arcs 
$\msf L^+_n X(R) = X(R[t]/t^n)$.
If $X$ is of finite type over $k$, then so is each $\msf L^+_n X$. If $X$ is smooth over $k$, then each $\msf L^+_n X$ is smooth over $k$
and $\msf L^+ X$ is formally smooth over $k$. If $X$ is affine, then so are
$\msf L^+_n X$ and $\msf L^+ X$. 

Define the formal loop space 
$\msf L X(R) = X(R\lbrac t)$. 
If $X$ is affine, then $\msf L X$ is representable by an ind-affine ind-scheme,
and we have a closed embedding $\msf L^+ X \into \msf L X$.

\subsubsection{}
Let $X$ be an affine spherical $G$-variety.
Define 
\[ \msf L^\bullet X := \msf L X \sm \msf L(X\sm X^\bullet), \]
which admits an open embedding into $\msf L X$. 
The $G$-action on $X$ induces a natural action of $\msf L^+ G$ on 
$\msf L X$ and $\msf L^+X, \msf L^\bullet X$ are stable under this action. 


\begin{eg}
\source{intersection-complexes-unramified-l-factors}
\notready \label{eg:A1}
\source{intersection-complexes-unramified-l-factors}
Let $X = \mbb A^1$ with the scaling $\mbb G_m$-action. Then  
$\msf L^+ X = \mbb A^\infty$, where we consider $\mbb A^\infty = \Spec k[a_0,a_1,\dotsc]$
as the coefficients of infinite Taylor series, and the ind-scheme
$\msf L X = ``\underset{\longrightarrow m}\lim" \Spec k[a_{-m},a_{-m+1},\dotsc]$ considered
as the coefficients of Laurent series.
Let $X^\bullet = \mbb A^1 \sm \{0\}$. Then $\msf L^\bullet X = \msf L X \sm \{0\}$ so $\msf L^+ X \cap \msf L^\bullet X = \mbb A^\infty \sm \{0\}$, whereas $\msf L^+(X^\bullet) = \mbb G_m \times \mbb A^\infty$. 
\end{eg}

\begin{rem}
\source{intersection-complexes-unramified-l-factors}
\notready 
The $k$-points of $\msf L^\bullet X$ are in bijection with $X^\bullet(k\lbrac t)$. 
Since $X^\bullet$ is in general not affine, however, $X^\bullet(k\lbrac t)$ does not
always have an ind-scheme structure. Even when $X^\bullet$ is affine, 
$\msf L(X^\bullet)$ may not be isomorphic to $\msf L^\bullet X$, despite having the same sets of $k$-points. 
\end{rem}


\subsubsection{Orbits on the formal loop space} \label{sect:stratLXpoint}
\source{intersection-complexes-unramified-l-factors}

For ease of notation, 
let $\mf o = k\tbrac t$ and $F = k\lbrac t$, so 
$\msf L^+G(k) = G(\mf o),\, \msf L^\bullet X(k) = X^\bullet(F)$.
We review the decomposition of $G(\mf o)$-orbits on $X^\bullet(F)$ due to \cite{LV}.
We present the reformulation of this result found in \cite[Theorem 3.3.1]{GN}. 

\smallskip

Let $X$ be an affine spherical variety, and pick a pair $(x_0, L)$ as in Theorem \ref{thm:localstructure} to write $X^\circ= T_X \times N_{P(X)}$ for its open Borel orbit. 
For a cocharacter $\ch\theta\in \ch\Lambda_X$ we let $t^{\ch\theta} \in T_X(F)$ 
denote the image of the uniformizer $t\in \mf o$ under the map $\ch\theta: \mbb G_m \to T_X$. 
Let $X^\bullet (F)/ G(\mf o)$ denote the \emph{set} of equivalence
classes of $G(\mf o)$-orbits in $X^\bullet (F)$.

\begin{thm}[{\cite{LV}, \cite[Theorem 3.3.1]{GN}}]
\source{intersection-complexes-unramified-l-factors}
\notready \label{thm:Gorbits}
\source{intersection-complexes-unramified-l-factors}
\begin{enumerate} 
\item The map 
\begin{equation}
 \ch\theta \mapsto 
X^\bullet(F)_{G:\ch\theta} := x_0 \cdot t^{\ch\theta} G(\mf o) 
\end{equation}
is a bijection of sets $\Cal V \cap \ch\Lambda_X \cong X^\bullet(F)/ G(\mf o)$ (which is independent of the choice of $(x_0,L)$). 
\item This bijection restricts to a bijection $\mf c_X^- 
\cong (X(\mf o)\cap X^\bullet(F))/ G(\mf o)$.
\end{enumerate}
\end{thm}

The theorem can be viewed as a generalization of the Cartan decomposition.
Define $\msf L^{\ch\theta} X$ to be the $\msf L^+G$-orbit of $x_0 t^{\ch\theta} \in \msf L^\bullet X$.



\begin{rem}
\source{intersection-complexes-unramified-l-factors}
\notready \label{rem:C0V}
\source{intersection-complexes-unramified-l-factors}
Observe that if $\ch\theta\in \mf c_X^-$, then 
in particular $\ch\theta\in \Lambda^-_X$ is antidominant (as a weight for the dual group $\ch G_X$ of $X$). 
Therefore, $w\ch\theta- \ch\theta\in \Lambda^\pos_X \subset \mf c_X$ for any $w\in W_X$, and hence $w\ch\theta\in \mf c_X$.
We suggest that the monoid $\mf c_X^- = \Cal C_0(X) \cap \Cal V \cap \ch\Lambda_X$ should be
thought of as the set of $W_X$-orbits in $\ch\Lambda_X$ that are entirely
contained in the cone $\Cal C_0(X)$.
\end{rem}


\begin{prop}
\source{intersection-complexes-unramified-l-factors}
\notready \label{prop:smoothfiberstrata}
\source{intersection-complexes-unramified-l-factors}
If $k$ has positive characteristic, 
assume that the stabilizer of $B$ acting on $x_0 \in X^\circ$ is a smooth subgroup. 
For any $\ch\theta \in \Cal V \cap \ch\Lambda_X$, 
the $\msf L^+G$-orbit $\msf L^{\ch\theta} X$ is a formally smooth $k$-scheme
with $k$-points equal to $X^\bullet(F)_{G:\ch\theta}$, and 
$\msf L^{\ch\theta} X$ is open in its closure in $\msf L X$. 
\end{prop}

Therefore, the collection of  
$\msf L^{\ch\theta}X,\, \ch\theta \in \Cal V \cap \ch\Lambda_X$, 
form a stratification of $\msf L^\bullet X$, in the sense that the strata are disjoint
and contain all the $k$-points.

\begin{proof}
Let $S \subset \msf L^+G$ denote the stabilizer subgroup of $x_0 t^{\ch\theta}
\in \msf L X(k)$. 
\emph{A priori}, the $\msf L^+G$-orbit $\msf L^{\ch\theta} X$ is 
defined as the fpqc sheaf quotient $\msf L^+G/S$. 
For $n \in \mbb N$, let $S_n$ denote the scheme-theoretic image 
of $S$ under the projection $\msf L^+G \to \msf L^+_n G$. 
For $m > n$, the transition map 
\[ \msf L^+_m G / S_m \to \msf L^+_n G / S_n = \msf L^+_m G / (S_n \cdot \ker(\msf L^+_m G \to \msf L^+_n G)) \]
is smooth affine since $\ker(\msf L^+_m G \to \msf L^+_n G)$ is smooth and unipotent. 
The schemes $\msf L^+_n G/S_n$ are also smooth since $\msf L^+_n G$ is smooth. 
Therefore $\msf L^{\ch\theta} X \cong \msf L^+ G / S \cong \varprojlim 
\msf L^+_n G / S_n$ is representable by a formally smooth scheme.
%Quotients of an irreducible solvable algebraic group over an algebraically closed field by an algebraic subgroup are always affine: Hochschild (1981, p. 184).

The fact that $\msf L^{\ch\theta} X \to \msf L X$ is open in its closure
will follow from Lemma~\ref{lem:theta-orbit-identification} below.
\end{proof}



\subsubsection{}
Recall the family $\mathscr X$ introduced in \S \ref{sect:affine-degeneration}. We define a map of formal loop spaces 
\begin{equation}\label{e:ithetaLX} 
\source{intersection-complexes-unramified-l-factors}
    i_{\ch\theta} : \msf L X \to \msf L \mathscr X 
\end{equation}
as follows:
By definition, we have a map $p: X \times T_X \to \mathscr X$. 
For a test ring $R$ and $\gamma \in X( R \lbrac t)$, 
define $i_{\ch\theta}(\gamma) = p( \gamma, t^{-\ch\theta} ) \in \mathscr X(R \lbrac t)$,
where we are considering $t^{-\ch\theta} \in T_X(R \lbrac t)$.

\begin{lem}
\source{intersection-complexes-unramified-l-factors}
\notready \label{lem:itheta-lands-open}
\source{intersection-complexes-unramified-l-factors}
The point $i_{\ch\theta}(x_0 t^{\ch\theta}) \in \msf L \mathscr X(k)$ is contained in 
$\msf L^+( \mathscr X^\circ )$. 
\end{lem}
\begin{proof}
By Proposition~\ref{prop:noncanon-section} (and the ensuing discussion), 
there exists a non-canonical section $s: T_X \times \ol{T_{X,\mss}} \to \mathscr X^\circ$
such that $s(1,1) = x_0$. 
If we choose a splitting $T \into G$ contained in the Levi $L$ from Theorem~\ref{thm:localstructure}, then the section $s$ is $T \times T_X$-equivariant, where 
$T_X$ acts diagonally on $T_X \times \ol{T_{X,\mss}}$. 
Thus on $k\lbrac t$-points, 
we have $i_{\ch\theta}(x_0 t^{\ch\theta}) = p(x_0 t^{\ch\theta}, t^{-\ch\theta})
= s(t^{\ch\theta}\cdot t^{-\ch\theta}, t^{-\ch\theta}) = s(1, t^{-\ch\theta})$,
which lies in $\mathscr X^\circ (k \tbrac t)$
since $-\ch\theta\in -\Cal V$. 
\end{proof}

The map $i_{\ch\theta}$ is $\msf L G$-equivariant. 
As a consequence of the lemma, we see that the $\msf L^+G$-orbit of
$i_{\ch\theta}(x_0 t^{\ch\theta})$ is contained in $\msf L^+(\mathscr X^\bullet)$. 
Thus
$i_{\ch\theta}$ restricts to a map $\msf L^{\ch\theta}X \to \msf L^+(\mathscr X^\bullet)$. 

We have a closed embedding $\msf L^+ \mathscr X \into \msf L \mathscr X$
and an open embedding $\msf L^+(\mathscr X^\bullet) \into \msf L^+\mathscr X$. 

\begin{lem}
\source{intersection-complexes-unramified-l-factors}
\notready  \label{lem:theta-orbit-identification}
\source{intersection-complexes-unramified-l-factors}
Assume that the stabilizer of $B$ acting on $x_0 \in X^\circ$ is a smooth subgroup. 
Then the map $i_{\ch\theta}$ induces an isomorphism 
\[ \msf L^{\ch\theta} X \overset\sim\to \msf L X \xt_{i_{\ch\theta}, \msf L \mathscr X} \msf L^+(\mathscr X^\bullet).
\]
\end{lem}
\begin{proof}
Since $i_{\ch\theta} : \msf L X \to \msf L \mathscr X$ is injective on $S$-points,
it suffices to show that 
$\msf L X \xt_{i_{\ch\theta}, \msf L \mathscr X} \msf L^+(\mathscr X^\bullet)$
is the $\msf L^+G$-orbit of $i_{\ch\theta}(x_0 t^{\ch\theta}) \in \msf L^+(\mathscr X^\bullet)(k)$.

Let $\Cal B = G/B$ denote the flag variety, and let
$(\mathscr X \times \Cal B)^\bullet \subset \mathscr X \times \Cal B$ denote 
the $G$-stable open subvariety 
consisting of all points $(x,\tilde B)$ where $x \in \mathscr X$ in the open $\tilde B$-orbit.
We have a smooth surjection $(\mathscr X\times \Cal B)^\bullet \to \mathscr X^\bullet$ by
first projection. 
On the other hand, Proposition~\ref{prop:noncanon-section} 
gives a non-canonical $G$-equivariant isomorphism 
$\mathscr X^\circ \xt^B G \cong (X^\circ \xt^B G) \times \ol{T_{X,\mss}}$. 
The action map and second projection induce an isomorphism 
$\mathscr X^\circ \xt^B G \cong (\mathscr X \times \Cal B)^\bullet$ and 
similarly for $X^\circ \xt^B G$. 
Thus there is a $G$-equivariant
isomorphism 
\[ (\mathscr X\times \Cal B)^\bullet \cong (X \times \Cal B)^\bullet \times \ol{T_{X,\mss}} \] 
over the base $\ol{T_{X,\mss}}$. 
The choice of base point $(x_0,B) \in (X \times \Cal B)^\bullet$ gives
an isomorphism $G/\textnormal{Stab}_B(x_0) \cong (X \times \Cal B)^\bullet$. 
Thus we have a $G$-equivariant smooth surjection 
$G \times \ol{T_{X,\mss}} \to \mathscr X^\bullet$ which induces
a surjection 
\begin{equation} \label{e:arc-family-orbit}
\source{intersection-complexes-unramified-l-factors}
    \msf L^+( G \times \ol{T_{X,\mss}} ) \to \msf L^+ (\mathscr X^\bullet)
\end{equation}
because $\on{Stab}_B(x_0)$ is assumed to be smooth. 
Since the composition 
$\msf L X \overset{i_{\ch\theta}}\to \msf L \mathscr X \to \msf L \ol{T_{X,\mss}}$ 
sends everything to the image of $t^{-\ch\theta}$ in $\msf L\ol{T_{X,\mss}}(k)$, 
we deduce that \eqref{e:arc-family-orbit} induces a 
surjection 
$\msf L^+G \to \msf L X \xt_{i_{\ch\theta}, \msf L \mathscr X} \msf L^+(\mathscr X^\bullet)$.
Thus the latter fiber product is a single $\msf L^+G$-orbit.
\end{proof}

\subsubsection{}
In this subsection we will use properties of toroidal compactifications of 
$X$. We refer the reader to \cite[\S 2.3-2.5]{SV}, \cite{KnLV}, \cite[\S 8]{GN} for an overview. 

Let $\ol X$ denote a complete, smooth toroidal embedding of $X^\bullet=H\bs G$ 
(the embedding is not 
simple if $N_G(H)/H$ is not finite). 
For any $\ch\theta \in \Cal V \cap \ch\Lambda_X$, the point 
$x_0 \cdot t^{\ch\theta} \in X^\bullet(F) \subset \ol X(F)$ defines
an $\mf o$-point of $\ol X$ by the valuative criterion of properness. 
In particular, we can take the limit as $t\to 0$ to get a point $x_{\ch\theta}
:= \lim_{t\to 0} x_0 \cdot t^{\ch\theta} \in\ol X(k)$.
Let $Z \subset \ol X$ denote the $G$-orbit of $x_{\ch\theta}$. 

Let $\Delta_X$ denote the set of normalized spherical roots of $X$.
Equivalently, $\Delta_X$ is the set of simple coroots of $\ch G_X$. 
Let $I \subset \Delta_X$ denote the set of spherical roots $\sigma$ 
such that $\brac{\sigma, \ch\theta}=0$. In the language of \cite[\S 2.3.6]{SV}, 
the orbit $Z$ ``belongs to $I$-infinity''. Let $X_I^\bullet = H_I\bs G$ denote 
the open $G$-orbit on the normal bundle $N_Z \ol X$ of $Z$ in $\ol X$; 
this is called a \emph{boundary degeneration} of $X$. 
As the notation suggests, $X_I^\bullet$ and the conjugacy class of $H_I$ depend only on the subset $I$
and not $\ch\theta$ (see \cite[Proposition 2.5.3]{SV}). We call $H_I$ a ``satellite'' of $H$,
following the terminology of \cite{BM}.


\begin{lem}
\source{intersection-complexes-unramified-l-factors}
\notready \label{lem:H_I-conn}
\source{intersection-complexes-unramified-l-factors}
Assume that $B$ acts simply transitively on $X^\circ$. 
Then the satellite subgroup $H_I \subset G$ is connected and smooth, for any subset $I\subset \Delta_X$.
\end{lem}
\begin{proof}
We have $X^\circ \cong X_I^\circ \cong B$. 
Let $H^0_I$ denote the reduced identity component of $H_J$.
Then $H^0_I\bs G \to H_I\bs G$ is a finite covering of spherical varieties sending the open $B$-orbit 
of $H^0_I\bs G$ to $X^\circ_I$. Since $B$ acts on $X^\circ_I\cong X^\circ$ with 
trivial stabilizer, the covering must be an isomorphism. 
\end{proof}

\begin{cor}
\source{intersection-complexes-unramified-l-factors}
\notready \label{cor:stab-conn}
\source{intersection-complexes-unramified-l-factors}
Assume that $B$ acts simply transitively on $X^\circ$. 
Let $\ch\theta \in \Cal V \cap \ch\Lambda_X$. 
Then the stabilizer of the group ind-scheme $\msf L H$ acting on $t^{\ch\theta} \in \Gr_G$
is a connected scheme.
\end{cor}
\begin{proof}
The cocharacter $-\ch\theta \in (-\Cal V)$ extends to a map 
$\mbb A^1 \to \ol{T_{X,\mss}}$. 
Let $\mf X := \mathscr X \xt_{\ol{T_{X,\mss}},-\ch\theta} \mbb A^1$. 
Then $\mf X \to \mbb A^1$ is an affine family with preimage over $\mbb G_m$ isomorphic
to $X \times \mbb G_m$. By \cite[Proposition 2.5.3]{SV}, the 
open $G$-orbit of the fiber over $0$ is isomorphic to $X_I^\bullet$. 
Let $\mf X^\bullet$ denote the union of the open $G$-orbits on each fiber,
so $\mf X^\bullet$ is a family over $\mbb A^1$ degenerating $X^\bullet$ to $X^\bullet_I$. 

The map $i_{\ch\theta}$ from \eqref{e:ithetaLX} 
factors through an embedding 
\[ \tilde i_{\ch\theta}: \msf L X \into \msf L \mf X \xt_{\msf L \mbb A^1} \pt \]
where $\pt \to \msf L \mbb A^1$ corresponds to $t \in k\lbrac t$. 
Lemma~\ref{lem:itheta-lands-open} implies that  
$\gamma := \tilde i_{\ch\theta}(x_0 t^{\ch\theta}) \in \msf L^+(\mf X^\bullet)$. 

Observe that $\on{Stab}_{\msf L H}(t^{\ch\theta}, \Gr_G) = 
\msf L H \cap t^{\ch\theta} (\msf L^+ G) t^{-\ch\theta}$ conjugates to 
\[ t^{-\ch\theta} (\msf L H) t^{\ch\theta} \cap \msf L^+G = 
\on{Stab}_{\msf L^+G}( x_0 \cdot t^{\ch\theta}, \msf L X ) = \on{Stab}_{\msf L^+G}( \gamma, \msf L^+(\mf X^\bullet) ). 
\]
Let $\mf J \subset G \times \mf X^\bullet$ denote the inertia group scheme
consisting of all $(g,x) \in G \times \mf X^\bullet$ such that $g x = x$. 
The fibers of $\mf J \to \mf X^\bullet$ are all conjugate to either $H$ or $H_I$. 
Since $\mf X^\bullet$ is smooth over $\mbb A^1$ (cf.~\cite[Corollary 5.2.3]{GN}), 
we have $H$ and $H_I$ are of the same dimension,
and they are smooth by Lemma~\ref{lem:H_I-conn}. 
Thus $\mf J \to \mf X^\bullet$ is a smooth morphism. 

Consider $\gamma$ as an arc $\Spec k\tbrac t \to \mf X^\bullet$. 
For $n\in \mbb N$, let $\gamma_n : \msf D_n := \Spec (k[t]/t^n) \to \mf X^\bullet$ denote the
corresponding $n$-jet and $\mf J_n := \mf J \xt_{\mf X^\bullet,\gamma_n} \msf D_n$ the
fiber product, which is a smooth group scheme over $\mf D_n$. 
Then 
\[ \on{Stab}_{\msf L^+G}( \gamma, \msf L^+(\mf X^\bullet) ) \cong 
\varprojlim_n \on{Sect}( \msf D_n, \mf J_n) 
\] 
where $\on{Sect}( \msf D_n, \mf J_n)$ 
is the scheme of maps $\msf D_n \to \mf J_n$ over $\msf D_n$. 
By smoothness of $\mf J \to \mf X^\bullet$, the
transition maps $\on{Sect}( \msf D_m, \mf J_m) \to \on{Sect}( \msf D_n, \mf J_n)$
are surjective for $m > n$. 
The transition maps are also group homomorphisms with unipotent kernels, so they
have contractible fibers. Therefore 
$\on{Stab}_{\msf L^+G}(\gamma, \msf L^+(\mf X^\bullet))$ contracts to 
$\on{Sect}(\msf D_0, \mf J_0) = \on{Stab}_G(\gamma(0), X_I^\bullet)$, which is conjugate
to the subgroup $H_I$. Since $H_I$ is connected by Lemma~\ref{lem:H_I-conn}, we are done.
\end{proof}
